\documentclass[11pt, a4paper]{article}

% Gemeinsame Style-Definitionen (TEXINPUTS muss templates/ enthalten — siehe scripts/build_tex.py)
% bfw-style.tex — gemeinsame Style-Definitionen für alle Handouts/Aufgabenhefte
% Voraussetzung: Kompilation mit xelatex (wegen fontspec)

\usepackage[ngerman]{babel}
% Babel-ngerman macht " zu einem aktiven Shorthand (z.B. für "a -> ä). In unseren
% Texten verwenden wir " als normales Zeichen — daher Shorthand komplett ausschalten.
\shorthandoff{"}
\usepackage{geometry}
\geometry{a4paper, top=2.4cm, bottom=2.4cm, left=2.3cm, right=2.3cm,
          headheight=15pt, headsep=0.7cm, footskip=1.1cm}

\usepackage{fontspec}
% Fallback-Kette: Source Sans Pro -> Calibri -> Segoe UI -> Latin Modern Sans
% (Latin Modern ist immer mit MiKTeX vorhanden)
\IfFontExistsTF{Source Sans Pro}{%
  \setmainfont{Source Sans Pro}[Ligatures=TeX]%
}{\IfFontExistsTF{Calibri}{%
  \setmainfont{Calibri}[Ligatures=TeX]%
}{\IfFontExistsTF{Segoe UI}{%
  \setmainfont{Segoe UI}[Ligatures=TeX]%
}{%
  \setmainfont{Latin Modern Sans}[Ligatures=TeX]%
}}}
\IfFontExistsTF{Source Code Pro}{%
  \setmonofont{Source Code Pro}[Scale=0.92]%
}{\IfFontExistsTF{Consolas}{%
  \setmonofont{Consolas}[Scale=0.92]%
}{%
  \setmonofont{Latin Modern Mono}[Scale=0.92]%
}}

% Gesperrte Versalien für Labels/Titelbalken (Kapitälchen-Ersatz, funktioniert
% mit jeder OpenType-Schrift — Latin Modern Sans hat keine echten Kapitälchen)
\newcommand{\bfwlabelfont}{\footnotesize\bfseries\addfontfeatures{LetterSpace=8.0}}

\usepackage{xcolor}
% Farbpalette: hell, freundlich, modern
\definecolor{bfwPrimary}{HTML}{0EA5A4}    % Petrol/Türkis - Primär
\definecolor{bfwPrimaryDark}{HTML}{0F766E} % dunkler Petrol für Headings
\definecolor{bfwAccent}{HTML}{F59E0B}     % Amber - Akzent
\definecolor{bfwSuccess}{HTML}{10B981}    % Grün - Erfolg / Lösung
\definecolor{bfwWarn}{HTML}{F97316}       % Orange - Achtung
\definecolor{bfwInk}{HTML}{0F172A}        % fast Schwarz
\definecolor{bfwInkSoft}{HTML}{475569}    % gedämpftes Grau
\definecolor{bfwBgSoft}{HTML}{F8FAFC}     % off-white
\definecolor{bfwBgTeal}{HTML}{ECFEFF}     % helles Türkis
\definecolor{bfwBgAmber}{HTML}{FEF3C7}    % helles Gelb
\definecolor{bfwBgRose}{HTML}{FFE4E6}     % helles Rosé für Eselsbrücken
\definecolor{bfwTabHead}{HTML}{E4F3F2}    % zarte Kopfzeilen-Hinterlegung für Tabellen

\usepackage[colorlinks=true, linkcolor=bfwPrimaryDark, urlcolor=bfwPrimaryDark]{hyperref}
\usepackage{lastpage}
\usepackage{enumitem}
\setlist[itemize]{leftmargin=*, itemsep=2pt, topsep=2pt}
\setlist[enumerate]{leftmargin=*, itemsep=2pt, topsep=2pt}

\usepackage[most]{tcolorbox}
\usepackage{booktabs}
\usepackage{colortbl}
\usepackage{tikz}
\usetikzlibrary{decorations.pathmorphing, positioning, shapes.geometric, arrows.meta}

% ---------------------------------------------------------------------------
% Einheitliches Tabellen-Design
%   - etwas mehr Zeilenluft, dezente graue Linien (wirkt auch mit \hline ruhiger)
%   - Kopfzeile einer Tabelle bei Bedarf zart hinterlegen: \rowcolor{bfwTabHead}
% ---------------------------------------------------------------------------
\renewcommand{\arraystretch}{1.25}
\arrayrulecolor{bfwInkSoft!50}
\setlength{\heavyrulewidth}{1pt}
\setlength{\lightrulewidth}{0.5pt}

% ---------------------------------------------------------------------------
% Boxen — klare farbige Titelbalken mit gesperrten Versal-Labels
% (Titel bleibt per Option überschreibbar: \begin{merksatz}[title={...}])
% ---------------------------------------------------------------------------
\tcbset{bfwbox/.style={%
  enhanced, breakable,
  boxrule=0.5pt, arc=2.5pt,
  left=10pt, right=10pt, top=8pt, bottom=8pt,
  toptitle=4pt, bottomtitle=4pt,
  titlerule=0pt,
  fonttitle=\bfwlabelfont,
  coltitle=white
}}

% Merkkästchen — wichtige Definition / Kernpunkt
\newtcolorbox{merksatz}[1][]{%
  bfwbox,
  colback=bfwBgTeal,
  colframe=bfwPrimaryDark,
  colbacktitle=bfwPrimaryDark,
  title={MERKE},
  #1
}

% Eselsbrücke — Mnemoniks
\newtcolorbox{eselsbruecke}[1][]{%
  bfwbox,
  colback=bfwBgRose,
  colframe=bfwAccent!85!black,
  colbacktitle=bfwAccent!85!black,
  title={ESELSBRÜCKE},
  #1
}

% Beispiel-Box
\newtcolorbox{beispiel}[1][]{%
  bfwbox,
  colback=bfwBgSoft,
  colframe=bfwInkSoft,
  colbacktitle=bfwInkSoft,
  title={BEISPIEL},
  #1
}

% Lösungs-Box (für Aufgabenhefte)
\newtcolorbox{loesung}[1][]{%
  bfwbox,
  colback=bfwSuccess!7,
  colframe=bfwSuccess!65!black,
  colbacktitle=bfwSuccess!65!black,
  title={LÖSUNG},
  #1
}

% Achtung-Box — typische Fehler / Stolperfallen
\newtcolorbox{achtung}[1][]{%
  bfwbox,
  colback=bfwWarn!6,
  colframe=bfwWarn!85!black,
  colbacktitle=bfwWarn!85!black,
  title={ACHTUNG},
  #1
}

% Formel-Box — für zentrale Formeln (groß, ruhig, ohne Titelbalken)
\newtcolorbox{formelbox}[1][]{%
  enhanced,
  colback=bfwBgTeal,
  colframe=bfwPrimaryDark,
  boxrule=1pt, arc=3pt,
  left=10pt, right=10pt, top=6pt, bottom=6pt,
  #1
}

% Glossar-Eintrag
\newcommand{\glossareintrag}[2]{%
  \par\noindent\textbf{\color{bfwPrimaryDark}#1} \enspace #2\par\smallskip
}

% ---------------------------------------------------------------------------
% Abschnitts-Design: farbiges Nummern-Quadrat + feine Linie unter \section,
% dezent farbige \subsection/\subsubsection
% ---------------------------------------------------------------------------
\usepackage{titlesec}
\newcommand{\bfwSecBadge}[1]{%
  \begin{tikzpicture}[baseline={([yshift=-0.62em]n.center)}]
    \node[fill=bfwPrimaryDark, text=white, font=\normalsize\bfseries,
          minimum width=1.55em, minimum height=1.55em, inner sep=1pt,
          rounded corners=1.5pt] (n) {#1};
  \end{tikzpicture}}
\titleformat{\section}
  {\Large\bfseries\color{bfwPrimaryDark}}
  {\bfwSecBadge{\thesection}}{0.6em}{}
  [\vspace{-0.2em}{\color{bfwPrimary!60}\titlerule[0.8pt]}]
\titleformat{\subsection}
  {\large\bfseries\color{bfwPrimaryDark}}
  {\textcolor{bfwPrimary}{\thesubsection}}{0.5em}{}
\titleformat{\subsubsection}
  {\normalsize\bfseries\color{bfwInk}}
  {\textcolor{bfwPrimary}{\thesubsubsection}}{0.4em}{}
\titlespacing*{\section}{0pt}{1.6em}{1em}
\titlespacing*{\subsection}{0pt}{1.2em}{0.5em}

% TikZ-Stil "skizzenhaft" — etwas weicher als Rough.js, aber stabil
\tikzset{
  roughbox/.style = {
    draw=bfwPrimaryDark, thick, fill=bfwBgTeal,
    rounded corners=4pt, inner sep=6pt
  },
  roughaccent/.style = {
    draw=bfwAccent, thick, fill=bfwBgAmber,
    rounded corners=4pt, inner sep=6pt
  },
  roughline/.style = {
    draw=bfwInkSoft, thick, -{Stealth[length=2.5mm]}
  }
}

% Bullet-Symbole (bewusst kein FontAwesome — vermeidet Font-Installations-Probleme)
\newcommand{\faLightbulb}{$\bullet$}
\newcommand{\faBookmark}{$\bullet$}

% ---------------------------------------------------------------------------
% Kopf-/Fußzeile
%   Kopf links:  Wortlogo „Wirtschaftsfachwirt"
%   Kopf rechts: Dokument-Kurztext, gesetzt via \kopfzeile{...}
%   Fuß links:   © Can Siebert · 2026 — Fuß rechts: Seite X von Y
%   Titelseite (Seite 1) bleibt ohne Kopf-/Fußzeile (\handouttitel setzt
%   \thispagestyle{empty}).
% ---------------------------------------------------------------------------
\usepackage{fancyhdr}
\newcommand{\bfwKopfzeileText}{}
\newcommand{\kopfzeile}[1]{\renewcommand{\bfwKopfzeileText}{#1}}
\pagestyle{fancy}
\fancyhf{}
\fancyhead[L]{\small\bfseries\color{bfwPrimaryDark}Wirtschaftsfachwirt}
\fancyhead[R]{\small\color{bfwInkSoft}\bfwKopfzeileText}
\renewcommand{\headrulewidth}{0.6pt}
\renewcommand{\headrule}{{\color{bfwPrimary}\hrule width\headwidth height 0.6pt}}
\fancyfoot[L]{\small\color{bfwInkSoft}\copyright\ Can Siebert\,\textperiodcentered\,2026}
\fancyfoot[R]{\small\color{bfwInkSoft}Seite \thepage\ von \pageref*{LastPage}}
% Auch \thispagestyle{plain} (z.B. durch \maketitle) soll leer bleiben:
\fancypagestyle{plain}{\fancyhf{}\renewcommand{\headrulewidth}{0pt}\renewcommand{\headrule}{}}

% ---------------------------------------------------------------------------
% \handouttitel{Obertitel}{Untertitel}{Kurs-Label}{Termin-Zeile}
%   Einheitlicher professioneller Titelblock: Dachzeile, farbige Headline,
%   Untertitel, farbige Trennlinie und dezente Meta-Zeile (Kurs/Termin/Dozent).
%   Ersetzt \title/\maketitle. Direkt nach \begin{document} aufrufen.
% ---------------------------------------------------------------------------
\newcommand{\handouttitel}[4]{%
  \thispagestyle{empty}%
  \begingroup
  \setlength{\parindent}{0pt}%
  {\bfwlabelfont\color{bfwPrimary}WIRTSCHAFTSFACHWIRT (IHK)\par}
  \vspace{0.55em}
  {\huge\bfseries\color{bfwPrimaryDark}#1\par}
  \vspace{0.5em}
  {\large\color{bfwInkSoft}#2\par}
  \vspace{0.9em}
  {\color{bfwPrimary}\rule{\linewidth}{2pt}}\par
  \vspace{0.55em}
  {\small\color{bfwInkSoft}#3\ \,\textperiodcentered\,\ #4\ \,\textperiodcentered\,\ Dozent: Can Siebert\par}
  \endgroup
  \vspace{1.4em}%
}


\kopfzeile{Handout VWL Lektion 3 \textperiodcentered{} Teil 2}

\begin{document}
\handouttitel{Lektion~3 \textperiodcentered{} Teil~2: Politikfelder, Außenwirtschaft \& EU}%
  {Von der Tarifautonomie bis zum Euro --- und die Wiederholung des Blocks 1.1}%
  {1.1 Volkswirtschaftliche Grundlagen}%
  {Sa, 29.08.2026 \textperiodcentered{} Session 2}

\begin{merksatz}[title={\bfseries Worum geht es in diesem Teil?}]
Nach Geld- und Finanzpolitik (Teil~1) lernen Sie in Teil~2 die weiteren
\textbf{wirtschaftspolitischen Maßnahmen} kennen: Wachstums-, Tarif-, Arbeitsmarkt- und
Umweltpolitik --- sowie die beiden großen \textbf{Konzeptionen} (nachfrage- vs.
angebotsorientierte Wirtschaftspolitik). Anschließend geht es um die
\textbf{Außenwirtschaft}: Freihandel und Protektionismus sowie die Besonderheiten der EU
(Binnenmarkt, Wirtschafts- und Währungsunion). Damit schließen wir den Themenblock~1.1
des Rahmenplans ab und wiederholen den gesamten Prüfungsstoff der Lektionen~1--3 ---
alle Themen sind regelmäßig Gegenstand der IHK-Prüfung.
\end{merksatz}

\tableofcontents
\vspace{1em}
\hrule
\vspace{1em}

\section{Wachstumspolitik}

Wachstumspolitik umfasst alle staatlichen Maßnahmen für ein \emph{angemessenes und
stetiges} Wirtschaftswachstum. Zu unterscheiden sind \textbf{quantitatives Wachstum}
(reine Zunahme des realen BIP bzw. Pro-Kopf-Einkommens) und \textbf{qualitatives Wachstum}
(zusätzlich Schonung der Umwelt und Verbesserung der Lebensqualität; Messkonzepte: soziale
Indikatoren der OECD, Net Economic Welfare). Ansatzpunkte der Wachstumspolitik:

\begin{itemize}
  \item Förderung von \textbf{Investitionen} (besseres Investitionsklima, Gewinnerwartung),
  \item Beschleunigung des technischen Fortschritts durch Förderung von
    \textbf{Forschung und Entwicklung},
  \item Förderung von \textbf{Bildung}, um das Angebot an qualifizierten Arbeitskräften zu
    erhöhen,
  \item im weiteren Sinne auch die \textbf{Wettbewerbspolitik} (funktionierender
    Leistungswettbewerb) und die wirtschaftliche Infrastruktur.
\end{itemize}

\section{Tarifpolitik}

Die Tarifpolitik umfasst alle Fragen rund um Verhandlung und Abschluss von
\textbf{Tarifverträgen} zwischen den Tarifparteien --- in Deutschland in erster Linie
\emph{Arbeitgeberverbände} und \emph{Gewerkschaften}; auch ein einzelner Arbeitgeber kann
Tarifvertragspartei sein (\emph{Haustarifvertrag}).

\begin{merksatz}[title={\bfseries Tarifautonomie (Art. 9 Abs. 3 GG)}]
Die Tarifpartner sind von Weisungen des Staates \textbf{unabhängig}. Der Abschluss von
Tarifverträgen liegt ausschließlich in den Händen der Tarifparteien; der Staat ist zur
Neutralität verpflichtet. Staatliche Vorgaben sind nur als \emph{Gesetze} zulässig
(z.\,B. Bundesurlaubsgesetz, Entgeltfortzahlungsgesetz, Arbeitszeitgesetz,
Mindestlohngesetz) --- sie setzen Mindestansprüche, die nicht unterschritten werden dürfen.
\end{merksatz}

\subsection{Tarifvertragsarten}

\begin{itemize}
  \item \textbf{Lohn- und Gehalts(rahmen)tarifverträge:} regeln die Entgeltfindung ---
    Höhe des Entgelts, Eingruppierung, Zuschläge, Zulagen;
  \item \textbf{Manteltarifverträge:} regeln die arbeitsrechtlichen Rahmenbedingungen ---
    Verteilung der Wochenarbeitszeit, Pausen, Urlaub, Schichtarbeit, Freistellungen,
    Altersversorgung (meist mit langer Laufzeit).
\end{itemize}

\subsection{Lohnpolitik: Kaufkrafttheorie vs. produktivitätsorientierte Lohnpolitik}

\begin{itemize}
  \item Die \textbf{Gewerkschaften} argumentieren \emph{nachfrageorientiert}
    (\textbf{Kaufkrafttheorie}): Höhere Löhne steigern die Nachfrage, regen Produktion und
    Investitionen an und verbessern die Beschäftigung.
  \item Die \textbf{Arbeitgeber} argumentieren \emph{angebotsorientiert}: Löhne sind ein
    Kostenfaktor, der die (internationale) Wettbewerbsfähigkeit gefährdet. Sie fordern eine
    \textbf{produktivitätsorientierte Lohnpolitik}: Lohnsteigerungen nicht höher als der
    Produktivitätsfortschritt --- das verhindert Lohnkosteninflation und lohnkostenbedingte
    Arbeitslosigkeit.
\end{itemize}

Lohnerhöhungen wirken doppelt: Sie erhöhen die Produktionskosten
(\emph{Kostenerhöhungseffekt}), können aber durch höhere Produktivität der eingesetzten
Arbeit kompensiert werden (\emph{Kosteneinsparungseffekt}). Volkswirtschaftlich schädlich
sind überhöhte Abschlüsse dann, wenn Unternehmen mit Preiserhöhungen, Rationalisierung oder
Produktionsverlagerung ins Ausland reagieren.

\section{Arbeitsmarktpolitik}

Der Arbeitsmarkt ist kein „freier Markt`` im Sinne der Preisbildung: Wegen der hohen
sozialen Bedeutung des Lohns besteht Konsens, dass der Arbeitsmarkt nicht vollständig dem
freien Spiel der Marktkräfte überlassen wird. Zuständig für die Arbeitsmarktpolitik im
engeren Sinne ist v.\,a. die \textbf{Bundesagentur für Arbeit} (Nürnberg) mit ihren
Regionalagenturen, Arbeitsagenturen und Jobcentern.

\begin{itemize}
  \item \textbf{Passive Arbeitsmarktpolitik:} mildert die wirtschaftlichen Folgen
    eingetretener oder drohender Arbeitslosigkeit --- \emph{Lohnersatzleistungen} wie
    Arbeitslosengeld I und II (Bürgergeld), Insolvenzgeld, Kurzarbeitergeld.
  \item \textbf{Aktive Arbeitsmarktpolitik:} versetzt Erwerbspersonen in die Lage, im
    Arbeitsprozess zu bleiben oder schnell wieder einzutreten --- Beratung und
    \emph{Vermittlung}, Aktivierung und berufliche Eingliederung (Trainings- und
    Qualifizierungsmaßnahmen), Förderung der \emph{Aus- und Weiterbildung}, Umschulung und
    berufliche Rehabilitation, \emph{Existenzgründungsförderung}, Mobilitätshilfen (Reise-
    und Umzugskosten).
\end{itemize}

Leitgedanke ist heute \textbf{„Fördern und Fordern``}: Arbeitslose müssen sich aktiv
bemühen und zumutbare Jobs annehmen, sonst drohen Kürzungen der Lohnersatzleistungen.
Kritikpunkte an der Arbeitsmarktpolitik: geschönte Arbeitslosenquoten (Geförderte fallen
aus der Statistik), teils ineffiziente Qualifizierungsmaßnahmen, Verschlechterung der
Qualität der Arbeitsverhältnisse (Befristung, Leiharbeit, Scheinselbstständigkeit).

\section{Umweltpolitik}

\subsection{Prinzipien der Umweltpolitik}

\begin{itemize}
  \item \textbf{Verursacherprinzip:} Die Kosten für Vermeidung, Beseitigung oder Ausgleich
    von Umweltbelastungen trägt der Verursacher --- externe Kosten werden
    \emph{internalisiert} (gehen in Kostenrechnung und Preise ein); Beispiele:
    Rücknahmepflicht für Umverpackungen, Abwasserabgabe.
  \item \textbf{Gemeinlastprinzip:} Ist kein Verursacher feststellbar, übernimmt der Staat
    die Maßnahmen --- finanziert über Steuern (z.\,B. kommunale Kläranlagen, Mülldeponien);
    Nachteil: keine Anreize für umweltbewusstes Handeln.
  \item \textbf{Vorsorgeprinzip:} Umweltgefahren vorausschauend vermeiden --- „besser
    Vorsorge statt Nachsorge`` (viele Schäden, z.\,B. das Ozonloch, sind kaum reparabel).
  \item \textbf{Kooperationsprinzip:} Zusammenarbeit von Bürgern, Unternehmen und Staat.
  \item \textbf{Nachhaltigkeitsprinzip} (sustainable development, ursprünglich aus der
    Forstwirtschaft): künftige Generationen nicht schlechterstellen --- erneuerbare
    Ressourcen nur im Umfang ihrer Regenerationsfähigkeit nutzen, nicht erneuerbare sparsam
    nutzen, Emissionen an der Aufnahmefähigkeit der Ökosysteme ausrichten (vgl. Agenda~21).
\end{itemize}

\subsection{Instrumente der staatlichen Umweltpolitik}

\begin{itemize}
  \item \textbf{Nicht fiskalische Instrumente:} \emph{Aufklärung und Information};
    \emph{Auflagen} (Ordnungsrecht): umweltbezogene Ge- und Verbote, z.\,B. Grenzwerte für
    Schadstoffausstoß, FCKW-Verbot, Umweltplaketten/Umweltzonen.
  \item \textbf{Fiskalische Instrumente:}
    \emph{Umweltabgaben} --- Steuern (Ökosteuer), Sonderabgaben (Abwasserabgabe), Gebühren
    (Abfall-/Abwassergebühren) verteuern umweltschädliches Verhalten;
    \emph{Umweltzertifikate} --- Umweltbelastung nur gegen Berechtigung: Ein Zertifikat
    gestattet eine bestimmte Menge Schadstoffausstoß pro Jahr; die Emissionsrechte werden
    börsenmäßig gehandelt (EU-Emissionshandel für CO\textsubscript{2});
    \emph{Umweltsubventionen} --- Zuschüsse oder Steuererleichterungen für
    umweltschützende Maßnahmen (z.\,B. Wärmedämmung).
\end{itemize}

\begin{eselsbruecke}
Die drei großen Instrumentenfamilien: \textbf{A-A-Z} --- \textbf{A}uflagen (Ordnungsrecht:
verbieten), \textbf{A}bgaben (verteuern), \textbf{Z}ertifikate (begrenzen und handeln).
\end{eselsbruecke}

\section{Nachfrage- vs. angebotsorientierte Wirtschaftspolitik}

Nachfrage- und Angebotstheorie liefern unterschiedliche idealtypische Konzepte zur
Bewältigung von Wachstums- und Beschäftigungskrisen. Sie unterscheiden sich darin, welcher
\emph{Marktseite} die größere Bedeutung beigemessen wird und welche \emph{Rolle der Staat}
spielen soll.

\subsection{Nachfragetheorie nach Keynes (Fiskalismus)}

Die Nachfragetheorie geht auf \textbf{John Maynard Keynes} (1883--1946) zurück und sieht
Wachstumsschwächen und Arbeitslosigkeit in einer \emph{unzureichenden gesamtwirtschaftlichen
Nachfrage} begründet. Da der Markt allein das Gleichgewicht nicht herstellt, soll der Staat
eine \textbf{aktive Rolle} übernehmen und für zusätzliche Nachfrage sorgen
(\emph{Symptombekämpfung}, kurzfristig): Stärkung der Massenkaufkraft durch Lohnerhöhungen
und staatliche Zuschüsse, Erhöhung des Staatskonsums durch öffentliche Ausgabenprogramme,
Ausweitung des öffentlichen Sektors --- notfalls über Verschuldung (\emph{Deficit
Spending}). Geldpolitisch soll die Zentralbank die Finanzpolitik unterstützen (Liquidität
bereitstellen, Zinsen niedrig halten); Fiskal- und Geldpolitik wirken \emph{antizyklisch}.
Anhänger heißen \textbf{Fiskalisten}. Grundlage der Lohnargumentation ist die
\emph{Kaufkrafttheorie}.

\textbf{Schwächen:} zeitliche Verzögerungen (\emph{Timelags}) zwischen Entscheidung und
Wirkung, Verdrängung privater Investitionen durch staatliche (\emph{Crowding-out}),
steigende Staatsquote sowie steigende Abgaben-, Steuer- und Schuldenlast.

\subsection{Angebotstheorie (Neoklassik, Monetarismus)}

Die Angebotstheorie geht auf die klassische Nationalökonomie zurück
(\textbf{Neoklassizismus}; prominenter Vertreter: \textbf{Milton Friedman}, 1912--2006).
Beschäftigung und Wachstum hängen demnach von den Bedingungen auf der
\emph{Angebotsseite} ab; der Staat soll sich \textbf{passiv} verhalten und nur die
\emph{Rahmenbedingungen} verbessern (\emph{Ursachenbekämpfung}, mittel- bis langfristig):
Kostendämpfung zur Erhöhung der Unternehmensrentabilität, Steuersenkungen, Abbau von
Subventionen und Regulierungen (Deregulierung), Verringerung des effizienzschwachen
Staatskonsums, Investitionsförderung, moderate Lohnabschlüsse. Motor der Entwicklung sind
die \emph{Investitionen der Unternehmen}: Bessere Renditeerwartungen führen zu
Investitionen und neuen Arbeitsplätzen. Die Forderung nach \emph{stetigem
Geldmengenwachstum} basiert auf dem \textbf{Monetarismus}: Zwischen Geldmenge und BIP
besteht eine enge Beziehung.

\textbf{Schwächen:} einseitige Betonung der Kosten --- Löhne sind auch ein
\emph{Nachfragefaktor}; eine bessere Gewinnsituation erhöht zwar die
Investitions\emph{fähigkeit}, nicht automatisch die Investitions\emph{bereitschaft} (bei
geringer Absatzerwartung bleiben Erweiterungsinvestitionen aus).

\begin{eselsbruecke}
\textbf{Keynes = Nachfrage = Staat aktiv = kurzfristig Symptome behandeln.}
\textbf{Angebot = Rahmenbedingungen = Staat passiv = langfristig Ursachen beheben.}
Merkhilfe: „\textbf{K}eynes \textbf{k}auft`` (Staat schafft Nachfrage) ---
„\textbf{F}riedman \textbf{f}reit die Märkte`` (Deregulierung).
\end{eselsbruecke}

\section{Außenwirtschaft: Freihandel oder Protektionismus}

Unter \textbf{Außenwirtschaft} versteht man die Gesamtheit aller wirtschaftlichen
Aktivitäten zwischen Wirtschaftssubjekten im Inland und im Ausland (Handel mit Gütern und
Dienstleistungen, Kapital- und Devisenverkehr; ausgewiesen in der \emph{Zahlungsbilanz}).
Gründe für internationalen Handel: unterschiedliche klimatische Bedingungen, ungleich
verteilte Rohstoff- und Energievorräte, unterschiedliche Spezialisierungen und
Technologieschwerpunkte --- Grundgedanke der \textbf{internationalen Arbeitsteilung}.

\subsection{Freihandel}

\textbf{Freihandel} ist der völlig unbehinderte internationale Güteraustausch --- der Staat
verzichtet auf marktbeeinflussende Maßnahmen. Vorteile: große Güterauswahl in verschiedenen
Qualitäten und Preislagen für die Verbraucher, \emph{Spezialisierungsvorteile} und
Kostensenkungen in der Produktion, Absatzmärkte im In- und Ausland, Innovationsförderung
durch internationalen Wettbewerb. Gewinner sind allerdings v.\,a. die starken,
international wettbewerbsfähigen Unternehmen.

\subsection{Protektionismus: tarifäre und nichttarifäre Handelshemmnisse}

Greift der Staat zum Schutz der einheimischen Wirtschaft in den freien Handel ein, spricht
man von \textbf{Protektionismus} (Motive: Schutz vor Billigimporten, Vermeidung von
Monostrukturen und Auslandsabhängigkeit).

\begin{itemize}
  \item \textbf{Tarifäre Handelshemmnisse} (preispolitisch): \emph{Importzölle} (größte
    Bedeutung --- verbessern die Wettbewerbsposition der einheimischen Wirtschaft und
    bringen Staatseinnahmen), Steuertarife, \emph{Exportsubventionen} (inländische Produkte
    können im Ausland günstiger angeboten werden). Innerhalb der EU werden keine Zölle
    erhoben; gegenüber Drittländern gilt ein einheitliches EU-Vorgehen (keine autonome
    deutsche Zollpolitik mehr).
  \item \textbf{Nichttarifäre Handelshemmnisse} (mengenpolitisch/administrativ):
    \emph{Importkontingente} (mengenmäßige Beschränkungen), \emph{Ein- und Ausfuhrverbote}
    (schärfste Form, z.\,B. Rüstungsgüter, Drogen), \emph{Embargos} (außen- oder
    sicherheitspolitisch angeordnete Beschränkungen, i.\,d.\,R. Umsetzung internationaler
    Sanktionen von UNO oder EU), administrative Erschwernisse (Abfertigungs- und
    Genehmigungsverfahren), technische Normen, Gesundheits-, Sicherheits- und
    Verbraucherschutzstandards --- letztere mit stark zunehmender Bedeutung.
\end{itemize}

\begin{merksatz}
Rechtsgrundlagen in Deutschland: \textbf{Außenwirtschaftsgesetz (AWG)} und
Außenwirtschaftsverordnung (AWV); das \textbf{BAFA} (Bundesamt für Wirtschaft und
Ausfuhrkontrolle) informiert über Verbote, Genehmigungen und Beschränkungen.
\textbf{Globalisierung} = weltweiter Austausch von Gütern, Dienstleistungen, Informationen
und Kapital sowie weltweite Vernetzung von Unternehmen. Die EU ist Mitglied der
\textbf{WTO} und schließt Freihandelsabkommen (z.\,B. CETA mit Kanada, JEFTA mit Japan).
\end{merksatz}

\section{Besonderheiten der EU}

Die Integration begann 1957 mit der \emph{Europäischen Wirtschaftsgemeinschaft (EWG)}
(Römische Verträge, sechs Gründungsstaaten); mit dem Vertrag von \textbf{Maastricht} (1993)
entstand die \emph{Europäische Union (EU)}.

\subsection{Der Europäische Binnenmarkt: die vier Freiheiten}

Ein \textbf{Binnenmarkt} ist ein Wirtschaftsraum ohne innere Grenzen (gegründet 1993). Er
ist durch \textbf{vier Freiheiten} gekennzeichnet:

\begin{center}
\begin{tabular}{p{4.2cm} p{9.6cm}}
\toprule
\textbf{Freiheit} & \textbf{Kennzeichen}\\
\midrule
Freier Personenverkehr & Abbau der Grenzkontrollen (Schengen), Aufenthalts- und Niederlassungsfreiheit, freie Arbeitsplatzwahl, Anerkennung von Berufsabschlüssen\\
Freier Warenverkehr & keine Zölle oder mengenmäßigen Beschränkungen, Harmonisierung von Umsatz-/Verbrauchssteuern und technischen Normen\\
Freier Dienstleistungsverkehr & Selbstständige und freie Berufe dürfen EU-weit anbieten und sich niederlassen; Liberalisierung bei Banken und Versicherungen\\
Freier Geld- und Kapitalverkehr & Öffnung der Finanzmärkte, erleichterter Zahlungsverkehr (SEPA), Aufhebung von Kapitalverkehrsbeschränkungen\\
\bottomrule
\end{tabular}
\end{center}

\textbf{Chancen:} größerer Absatzmarkt, sinkende Stückkosten durch höhere Stückzahlen und
Rationalisierung, vielfältigeres und günstigeres Güterangebot, stabilisierende Wirkung auf
das Preisniveau. \textbf{Risiken/Nachteile:} Verlagerung von Produktionsstätten ins
kostengünstigere EU-Ausland (Arbeitsplatzverluste), Bürokratiekosten, aus Sicht der
Kritiker „Entmachtung`` der Einzelstaaten.

\begin{eselsbruecke}
Die vier Freiheiten: \textbf{P-W-D-K} --- \textbf{P}ersonen, \textbf{W}aren,
\textbf{D}ienstleistungen, \textbf{K}apital. Merkspruch: „\emph{Praktisch wandern durch
Kontinente}``.
\end{eselsbruecke}

\subsection{Europäische Wirtschafts- und Währungsunion (EWWU)}

Die Europäische Währungsunion ist Teil der EWWU; die Teilnehmerstaaten bilden die
\textbf{Eurozone}. Nach dem Maastrichter Vertrag dürfen nur EU-Staaten den Euro einführen,
die ihre Stabilität durch die \textbf{Konvergenzkriterien} bewiesen haben:

\begin{enumerate}
  \item \textbf{Inflationsrate:} im Jahr vor dem Beitritt höchstens 1,5~Prozentpunkte über
    dem Mittelwert der drei preisstabilsten EU-Staaten;
  \item \textbf{Nettoneuverschuldung:} nicht dauerhaft mehr als \textbf{3~\%} des BIP;
  \item \textbf{Gesamtschuldenstand:} i.\,d.\,R. nicht mehr als \textbf{60~\%} des BIP;
  \item \textbf{Zinssatz für langfristige Kredite:} höchstens 2~Prozentpunkte über dem
    Mittelwert der drei preisstabilsten EU-Staaten;
  \item \textbf{Wechselkursstabilität:} zwei Jahre stabiler Außenwert ohne Abwertung
    (Teilnahme am Wechselkursmechanismus WKM~II).
\end{enumerate}

Die Einhaltung wurde in der Vergangenheit nicht konsequent überprüft; mit der Euro- und
Staatsschuldenkrise reformierten EU-Parlament und Finanzminister 2011 den europäischen
\emph{Stabilitäts- und Wachstumspakt} (härteres Vorgehen gegen „Haushaltssünder``; die
Referenzwerte 3~\% und 60~\% blieben erhalten).

\subsection{Vor- und Nachteile der gemeinsamen Währung}

\begin{itemize}
  \item \textbf{Vorteile:} keine Wechselkursrisiken und Umtauschkosten im Euroraum,
    Preistransparenz (direkter Preisvergleich), Förderung von Handel, Investitionen und
    Reiseverkehr, starke Währung im Weltmaßstab.
  \item \textbf{Nachteile:} Verzicht auf eigenständige nationale Geld- und
    Wechselkurspolitik (keine Abwertung zur Stärkung der Wettbewerbsfähigkeit möglich),
    einheitlicher Leitzins für unterschiedlich entwickelte Volkswirtschaften,
    Haftungs- und Transferrisiken bei Schuldenkrisen einzelner Mitglieder.
\end{itemize}

\section{Formeln \& Rechenbeispiele}

Die folgenden Berechnungen sind Prüfungsklassiker --- alle folgen demselben Muster:
\emph{Teil durch Ganzes mal 100}. Prägen Sie sich Formel und Rechenweg anhand der
einfachen Zahlenbeispiele ein.

\begin{merksatz}[title={\bfseries Formel: Arbeitslosenquote}]
\[
\text{Arbeitslosenquote in \%} =
\frac{\text{registrierte Arbeitslose}}{\text{alle zivilen Erwerbspersonen}} \times 100
\]
Zivile Erwerbspersonen = Erwerbstätige \emph{plus} registrierte Arbeitslose (also alle,
die arbeiten oder arbeiten wollen).
\end{merksatz}

\begin{beispiel}
In einem Land sind 2,3~Mio. Personen arbeitslos gemeldet; es gibt 43,7~Mio. Erwerbstätige.\\[0.3em]
\textbf{Schritt 1:} Erwerbspersonen bestimmen: $43{,}7 + 2{,}3 = 46$~Mio.\\
\textbf{Schritt 2:} Arbeitslose durch Erwerbspersonen teilen: $2{,}3 \div 46 = 0{,}05$.\\
\textbf{Schritt 3:} Mal 100: $0{,}05 \times 100 = \textbf{5{,}0~\%}$ Arbeitslosenquote.
\end{beispiel}

\begin{merksatz}[title={\bfseries Anwendung: Konvergenzkriterien nachrechnen (3~\% / 60~\%)}]
\[
\text{Neuverschuldungsquote} = \frac{\text{jährliche Neuverschuldung}}{\text{BIP}} \times 100
\qquad
\text{Schuldenstandsquote} = \frac{\text{Gesamtschulden}}{\text{BIP}} \times 100
\]
Grenzwerte: Neuverschuldung max. \textbf{3~\%}, Gesamtschuldenstand max. \textbf{60~\%}
des BIP.
\end{merksatz}

\begin{beispiel}
Ein Beitrittskandidat hat ein BIP von 500~Mrd.~Euro, eine Neuverschuldung von
20~Mrd.~Euro und Gesamtschulden von 280~Mrd.~Euro.\\[0.3em]
\textbf{Schritt 1:} Neuverschuldung: $20 \div 500 \times 100 = 4{,}0~\%$ ---
\emph{über} der 3-\%-Grenze~$\Rightarrow$ Kriterium \textbf{verletzt}.\\
\textbf{Schritt 2:} Schuldenstand: $280 \div 500 \times 100 = 56{,}0~\%$ ---
\emph{unter} der 60-\%-Grenze~$\Rightarrow$ Kriterium \textbf{erfüllt}.\\
\textbf{Ergebnis:} Das Land erfüllt die fiskalischen Konvergenzkriterien insgesamt
\emph{nicht} (beide Grenzwerte müssen eingehalten werden).
\end{beispiel}

\begin{merksatz}[title={\bfseries Wichtig für die Prüfungsvorbereitung (Block 1.1 komplett)}]
Sie sollten sicher können: die Gegenüberstellung \emph{nachfrage- vs. angebotsorientiert}
(Keynes vs. Friedman), \emph{tarifäre vs. nichttarifäre Handelshemmnisse}, die \emph{vier
Freiheiten} des Binnenmarkts, die \emph{Konvergenzkriterien} (3~\%~/ 60~\%), die
Prinzipien und Instrumente der \emph{Umweltpolitik} (A-A-Z), \emph{aktive vs. passive
Arbeitsmarktpolitik}, die \emph{Tarifautonomie} mit den Tarifvertragsarten sowie die
\emph{Formeln} für Arbeitslosenquote und Konvergenzquoten (Abschnitt~8). Wiederholen Sie
außerdem Teil~1 dieser Lektion (Konjunktur, EZB, Fiskalpolitik),
Lektion~1 (Preisbildung, Marktformen) und Lektion~2 (Wettbewerb, Staatseingriffe, VGR,
magisches Viereck/Sechseck) --- die Quiz beider Teile und das Prüfungstraining helfen dabei.
\end{merksatz}

\vspace{1em}
\hrule
\vspace{0.8em}
\noindent\textsf{\small Quellen: Rahmenplan Wirtschaftsbezogene Qualifikationen (DIHK), Abschnitte 1.1.3.2 und 1.1.4;
Textband Volkswirtschaft (DIHK-Bildungs-GmbH), Kap.~3.2 und Kap.~4, S.~45--64. Nächster Termin (Sa, 05.09.2026):
Start des BWL-Teils --- Betriebliche Funktionen~I.}

\end{document}
